% !Mode:: "TeX:UTF-8"
% cSpell:ignore threshv topk randomk warmup PPL tok

\chapter{混合压缩与自适应调度的设计与实现}

第3章实现的9种通信压缩算法可分为量化与稀疏化两大类别，二者分别从``降低比特宽度''和``减少传输元素''两个方向压缩通信量。然而，单一方向的压缩手段在实践中各有局限：量化方法在保持可接受收敛性的条件下，压缩比通常不超过32倍；稀疏化方法虽然名义压缩率很高，但不可避免地要在通信消息中携带索引，使得实际带宽节省被部分抵消。与此同时，第4章的对比实验揭示了一个值得关注的现象——训练不同阶段对压缩强度的耐受程度差异明显：早期梯度幅值大、方向变化剧烈，对压缩噪声的容忍度相对更高；后期梯度趋于平缓后，同等绝对误差所占的相对比例增大，对收敛精度的影响也就更大。

基于以上分析，本章从两个方向对通信压缩进行扩展：（1）\textbf{两阶段混合压缩}，将稀疏化与量化级联以获得乘性压缩比；（2）\textbf{自适应压缩调度}，根据训练进度或梯度统计量动态调整压缩强度。两者都在第3章的FSDP通信钩子框架下实现，可以直接与已有算法配合使用。

\section{两阶段混合压缩}

\subsection{设计动机}

Top-k稀疏化（\(\rho=0.01\)）将参与通信的梯度元素数量从\(n\)降至\(k = \rho \cdot n\)，int8量化则将每个梯度元素的位宽从32比特压到8比特。如果将两者级联——先稀疏化筛选出少量关键梯度元素，再对这些元素做低比特量化——压缩效果就具有乘性关系：
\begin{equation}
r_{\text{total}} = r_{\text{sparse}} \times r_{\text{quant}} = \frac{1}{\rho} \times \frac{32}{8} = 100 \times 4 = 400.
\label{eq:hybrid-ratio}
\end{equation}
其中\(r_{\text{sparse}}\)为稀疏化压缩比，\(r_{\text{quant}}\)为量化压缩比。若将第二阶段从int8换成1-bit编码（仅保留符号位），\(r_{\text{quant}} = 32\)，理论总压缩比可达\(100 \times 32 = 3200\)倍。

当然，实际通信中还需传输每个选中梯度元素的索引（int64, 8字节），因此单rank实际发送字节数为：
\begin{equation}
B_{\text{hybrid}} = k \times (b_{\text{index}} + b_{\text{quant}}) + b_{\text{meta}},
\label{eq:hybrid-bytes}
\end{equation}
其中\(b_{\text{index}}=8\)字节为int64索引开销，\(b_{\text{quant}}\)为量化后单元素字节数（int8为1字节，1-bit为\(1/8\)字节），\(b_{\text{meta}}\)为缩放因子或均值参数等少量元数据（通常\(\le 8\)字节）。对比无压缩的\(B_{\text{baseline}} = n \times 4\)字节（float32），有效压缩比为：
\begin{equation}
r_{\text{effective}} = \frac{n \times 4}{k \times (8 + b_{\text{quant}}) + b_{\text{meta}}}.
\label{eq:hybrid-effective}
\end{equation}

以\(n = 10^6\)、\(\rho = 0.01\)为例，Top-k+int8的有效压缩比为\(4 \times 10^6 / (10^4 \times 9 + 4) \approx 44.4\)倍，Top-k+1-bit约49.2倍。索引开销使有效压缩比低于理论乘性值，但相比单一方法仍有数量级提升。

\subsection{通用混合压缩框架}

为支持稀疏化与量化方法的灵活组合，本文设计了一个两阶段混合压缩框架，其核心流程可抽象为：
\begin{equation}
\mathbf{c} = Q(S(\tilde{\mathbf{g}})), \quad \hat{\mathbf{g}} = D_Q(D_S(\mathbf{c})),
\label{eq:hybrid-abstract}
\end{equation}
其中\(S(\cdot)\)为稀疏化算子，\(Q(\cdot)\)为量化算子，\(D_S(\cdot)\)和\(D_Q(\cdot)\)分别为对应的解码算子。\(\tilde{\mathbf{g}} = \mathbf{g}_t + \mathbf{e}_t\)为叠加了误差反馈残差的梯度。

图~\ref{fig:hybrid-pipeline}展示了以Top-k+int8为例的完整两阶段混合压缩流程。其输入与第3章的总流程图保持一致：先从完整梯度\(\mathbf{g}_{\text{full}}\)出发，与残差张量叠加后得到\(\tilde{\mathbf{g}}\)。在第一阶段，Top-k保留绝对值最大的\(k\)个梯度元素，得到索引集合\(\mathcal{I}\)与对应取值向量\(\mathbf{v}\)；在第二阶段，\(\mathbf{v}\)经过int8量化得到紧凑表示\(\mathbf{c}\)。压缩后的索引、量化值和元数据通过all\_gather在所有rank间交换，接收端逐rank解码后将各rank的贡献叠加，提取本rank分片并除以world size，最终输出平均梯度分片\(\mathbf{g}_{\text{shard}}\)。最后，两阶段的总误差被计算并累积到残差张量中，用于下一步的误差反馈。

\begin{figure}[htbp]
\centering
\begin{tikzpicture}[
    box/.style={draw, rounded corners, minimum width=2.4cm, minimum height=0.85cm, align=center, font=\small},
    smallbox/.style={draw, rounded corners, minimum width=2.0cm, minimum height=0.7cm, align=center, font=\footnotesize},
    arrow/.style={->, thick, >=stealth}
]
\node[box, fill=blue!10] (input) at (0,0) {完整梯度\\$\mathbf{g}_{\text{full}}$};
\node[box, fill=orange!15] (efadd) at (3.1,0) {误差反馈叠加\\$\tilde{\mathbf{g}}$};
\node[box, fill=green!15] (sparse) at (6.2,0) {Stage 1\\稀疏化};
\node[box, fill=green!15] (quant) at (9.3,0) {Stage 2\\量化};
\node[box, fill=yellow!15] (comm) at (9.3,-2.0) {跨rank交换梯度};
\node[box, fill=green!15] (decode) at (6.2,-2.0) {解码还原\\$D_Q\circ D_S$};
\node[box, fill=blue!10] (output) at (3.1,-2.0) {平均梯度分片\\$\mathbf{g}_{\text{shard}}$};

\draw[arrow] (input) -- (efadd);
\draw[arrow] (efadd) -- (sparse);
\draw[arrow] (sparse) -- (quant);
\draw[arrow] (quant) -- (comm);
\draw[arrow] (comm) -- (decode);
\draw[arrow] (decode) -- (output);

\node[smallbox, fill=orange!10, minimum width=3.3cm] (idx) at (5.9,1.6) {输出：索引集合\(\mathcal{I}\)\\对应取值向量\(\mathbf{v}\)};
\node[smallbox, fill=orange!10, minimum width=2.8cm] (qval) at (9.6,1.6) {输出：量化压缩\\\(\mathbf{c}\)};
\draw[arrow, dashed, gray] (sparse.north) -- (idx.south);
\draw[arrow, dashed, gray] (quant.north) -- (qval.south);

\node[smallbox, fill=red!10, minimum width=5.5cm] (ef) at (6.2,-3.8) {残差更新：$\mathbf{e}_{t+1}=\tilde{\mathbf{g}}-D_Q(D_S(\mathbf{c}))$};
\draw[arrow, dashed] (decode.south) -- (ef.north);

\node[font=\footnotesize] at (6.2,2.35) {Top-k / Threshold-v / Random-k};
\node[font=\footnotesize] at (9.3,2.35) {int8 / 1-bit};
\end{tikzpicture}
\caption{两阶段混合压缩流程（以Top-k+int8为例）}
\label{fig:hybrid-pipeline}
\end{figure}

框架通过HybridCompressState状态类统一管理稀疏化方法（sparse\_method）、量化方法（quant\_method）、稀疏率（ratio）等超参数。当前实现支持表~\ref{tab:hybrid-combos}中的5种组合，其中第一列为对应的注册名称。

\begin{table}[htbp]
\caption{两阶段混合压缩的组合方案}
\label{tab:hybrid-combos}
\vspace{0.5em}\centering\wuhao
\begin{tabular}{llccc}
\toprule
注册名称 & Stage-1 \(\to\) Stage-2 & 稀疏率\(\rho\) & 每元素传输位宽 & 理论乘性压缩比 \\
\midrule
topk\_int8     & Top-k \(\to\) int8   & 0.01 & \(8+64\) bits & \(100 \times 4\) \\
topk\_1bit     & Top-k \(\to\) 1-bit  & 0.01 & \(1+64\) bits & \(100 \times 32\) \\
thresholdv\_int8 & Threshold-v \(\to\) int8 & 0.01 & \(8+64\) bits & \(100 \times 4\) \\
thresholdv\_1bit & Threshold-v \(\to\) 1-bit & 0.01 & \(1+64\) bits & \(100 \times 32\) \\
randomk\_int8  & Random-k \(\to\) int8 & 0.01 & \(8+64\) bits & \(100 \times 4\) \\
\bottomrule
\end{tabular}
\end{table}

\subsection{第一阶段：稀疏化}

第一阶段从梯度\(\tilde{\mathbf{g}} \in \mathbb{R}^n\)中选取\(k\)个梯度元素，得到索引集合\(\mathcal{I} \subseteq \{0, 1, \ldots, n-1\}\)及对应取值向量\(\mathbf{v} \in \mathbb{R}^{|\mathcal{I}|}\)。三种可用方法在选取策略上各有侧重。

Top-k选取按绝对值降序保留前\(k\)个梯度元素。这一操作是有偏的，因此需配合误差反馈来补偿被丢弃的信息。

Threshold-v选取保留绝对值超过阈值\(v\)的梯度元素。未显式指定阈值时，系统利用Top-k操作估计使恰好保留\(k\)个梯度元素的隐式阈值，并对超出的候选做二次筛选以严格控制通信量上界。

Random-k选取从\([n]\)中均匀采样\(k\)个元素索引，并将选中的梯度元素乘以\(n/k\)以保持无偏性。由于选取不依赖梯度值分布，计算开销在三者中最低。

\subsection{第二阶段：量化}

在稀疏化选出的\(|\mathcal{I}|\)个梯度元素上，进一步做低比特量化以压缩每个元素的表示开销。

int8对称量化沿用第3章的方案，对对应取值向量\(\mathbf{v}\)计算局部缩放因子\(s = q_{\max} / (\max|\mathbf{v}| + \epsilon)\)并量化为int8整数。每个选中梯度元素传输1字节量化值和8字节索引，外加1个float32缩放因子。

1-bit符号量化仅保留符号\(\text{sign}(\mathbf{v}_i)\)，并分别计算正值元素和负值元素的均值作为重构参考：
\begin{equation}
\mu^+ = \frac{1}{|\{i : \mathbf{v}_i > 0\}|} \sum_{i: \mathbf{v}_i > 0} \mathbf{v}_i, \quad
\mu^- = \frac{1}{|\{i : \mathbf{v}_i < 0\}|} \sum_{i: \mathbf{v}_i < 0} \mathbf{v}_i.
\label{eq:hybrid-1bit-mu}
\end{equation}
解码时根据符号选择\(\mu^+\)或\(\mu^-\)进行重构。每个梯度元素传输1字节符号和8字节索引，附带2个float32均值参数。

\subsection{通信模式}

压缩后的紧凑表示通过all\_gather\_into\_tensor在所有rank间交换。每个rank发送三份数据：索引张量（int64）、量化值张量（int8或int8符号）以及元数据张量（缩放因子或均值对）。接收端逐rank解码后，通过index\_add\_将各rank的贡献叠加到全零张量上，提取本rank分片并除以world size得到平均梯度分片。

值得一提的是，为保证all\_gather输入在各rank间维度一致，当Threshold-v实际选中的梯度元素少于\(k\)时，需对索引和量化值做零填充；解码时填充位置的零值不会影响index\_add\_的正确性。

\subsection{误差反馈}

两阶段级联带来叠加的近似误差。设稀疏化误差\(\mathbf{e}_{\text{sparse}} = \tilde{\mathbf{g}} - S(\tilde{\mathbf{g}})\)，量化误差\(\mathbf{e}_{\text{quant}} = S(\tilde{\mathbf{g}}) - D_Q(Q(S(\tilde{\mathbf{g}})))\)，总误差为二者之和：
\begin{equation}
\mathbf{e}_{\text{total}} = \mathbf{e}_{\text{sparse}} + \mathbf{e}_{\text{quant}} = \tilde{\mathbf{g}} - D_Q(Q(S(\tilde{\mathbf{g}}))).
\label{eq:hybrid-error}
\end{equation}
误差反馈机制将\(\mathbf{e}_{\text{total}}\)作为下一步的残差\(\mathbf{e}_{t+1}\)，使被丢弃的信息在后续迭代中得以补偿。这与第3章单一压缩方法的误差反馈公式形式一致，区别仅在于压缩算子\(C(\cdot)\)变成了\(Q \circ S\)的级联。

\section{自适应压缩调度}

\subsection{设计动机}

第4章的实验结果中，有一个尚未被充分利用的观察：训练初期梯度幅值大且变化剧烈，压缩引入的相对误差对优化方向的干扰有限；而随着训练推进，梯度逐渐平缓，相同绝对误差的相对占比反而增大，对收敛精度的影响也就更明显。这种``不同阶段对压缩噪声的敏感程度不同''的特点，提示了一种在时间维度上动态调节压缩强度的可能性——在不同阶段施加不同程度的压缩，从而更好地平衡通信节省与收敛质量。

这一思路可以形式化为将稀疏率\(\rho\)从固定常数推广为训练步数\(t\)或梯度统计量的函数：
\begin{equation}
\rho(t): \{0, 1, \ldots, T-1\} \to [\rho_{\min}, \rho_{\max}],
\label{eq:adaptive-schedule}
\end{equation}
其中\(T\)为总训练步数，\(\rho_{\min}\)和\(\rho_{\max}\)分别为稀疏率的下限和上限。\(\rho\)越小表示压缩越激进（保留的梯度元素越少），\(\rho\)越大则压缩越温和。

\subsection{调度策略}

本文实现了三种调度策略，分别面向不同的训练场景和调试需求。

\subsubsection{线性调度}

线性调度提供最直接的基线方案：\(\rho\)从\(\rho_{\max}\)均匀下降至\(\rho_{\min}\)。
\begin{equation}
\rho_{\text{linear}}(t) = \rho_{\max} - \frac{t}{T}(\rho_{\max} - \rho_{\min}).
\label{eq:linear-schedule}
\end{equation}
实现简单、行为可预测，主要用作消融实验中的对照组。

\subsubsection{Warmup-Decay调度}

Warmup-Decay将训练过程分为预热与衰减两个阶段。预热阶段（\(0 \le t < T_w\)）内，稀疏率从\(\rho_{\max}\)线性下降至\(\rho_{\min}\)：
\begin{equation}
\rho_{\text{warmup}}(t) = \rho_{\max} - \frac{t}{T_w}(\rho_{\max} - \rho_{\min}), \quad t < T_w,
\label{eq:warmup}
\end{equation}
其中\(T_w = \lfloor \alpha \cdot T \rfloor\)为预热步数，\(\alpha\)为预热比例（默认取0.1）。其背后的直觉是：训练初期模型参数远离最优解，梯度方向的方差虽大但幅值也大，保留更多元素有助于快速下降。

衰减阶段（\(T_w \le t < T\)）内，稀疏率按余弦曲线从\(\rho_{\min}\)逐步回升至\(\rho_{\text{mid}} = (\rho_{\min} + \rho_{\max})/2\)：
\begin{equation}
\rho_{\text{decay}}(t) = \rho_{\min} + (\rho_{\text{mid}} - \rho_{\min}) \cdot \left(1 - \frac{1}{2}\left(1 + \cos\frac{\pi(t - T_w)}{T - T_w}\right)\right), \quad t \ge T_w.
\label{eq:decay}
\end{equation}
训练中期可以容忍最强压缩以最大化通信节省，而尾部适当放松则有助于微调阶段的精度收敛。

\subsubsection{梯度自适应调度}

前两种策略均为基于步数的确定性函数，不感知当前梯度的实际统计特征。梯度自适应调度则根据实时梯度范数动态决策。维护梯度\(\ell_2\)范数的指数移动平均（EMA）：
\begin{equation}
\bar{g}_t = \beta \cdot \|\mathbf{g}_t\|_2 + (1-\beta) \cdot \bar{g}_{t-1},
\label{eq:grad-ema}
\end{equation}
其中\(\beta = 0.1\)，\(\bar{g}_0 = 0\)。当\(\|\mathbf{g}_t\|_2\)显著大于\(\bar{g}_t\)时，说明当前步梯度携带较强信号，应减轻压缩；反之则可加大压缩力度。

稀疏率通过sigmoid函数从范数比值映射到\([\rho_{\min}, \rho_{\max}]\)区间：
\begin{equation}
\rho_{\text{grad}}(t) = \rho_{\min} + \frac{\rho_{\max} - \rho_{\min}}{1 + \exp\left(-2 \cdot \left(\frac{\|\mathbf{g}_t\|_2}{\bar{g}_t} - 1\right)\right)}.
\label{eq:grad-adaptive}
\end{equation}

当\(\|\mathbf{g}_t\|_2 / \bar{g}_t \gg 1\)时sigmoid输出趋近1，\(\rho \to \rho_{\max}\)（轻压缩）；比值远小于1时输出趋近0，\(\rho \to \rho_{\min}\)（强压缩）。这使得压缩策略能够自动跟踪训练过程中梯度信号强度的变化，而不需要预设与步数相关的超参数。

图~\ref{fig:schedule-curves}对比了三种调度策略下稀疏率\(\rho\)随训练步数\(t\)的变化曲线。设定\(\rho_{\min} = 0.001\)、\(\rho_{\max} = 0.10\)、\(T = 200\)、\(\alpha = 0.1\)。可以看到，线性调度从\(\rho_{\max}\)单调下降至\(\rho_{\min}\)，行为最为简单；Warmup-Decay在预热阶段从\(\rho_{\max}\)快速下降到\(\rho_{\min}\)，随后在衰减阶段按余弦曲线回升到\(\rho_{\text{mid}}\)，整体呈现``先紧后松''的形态；梯度自适应调度的曲线并非固定形状，而是随着实际梯度范数的变化而波动，图中给出的是一次典型运行的轨迹，其在训练前期与Warmup-Decay的行为有相似之处，但在训练后期会根据梯度信号强度自适应调节。

\begin{figure}[htbp]
\centering
\includegraphics[width=0.82\textwidth]{fig9}
\caption{三种自适应调度策略的稀疏率变化曲线对比}
\label{fig:schedule-curves}
\end{figure}

在实现层面，自适应调度器作为外层包装，内部调用稀疏化方法（Top-k、Random-k或Threshold-v）完成梯度元素筛选，并默认使用int8作为第二阶段量化。每次hook调用时，调度器先根据当前步数或梯度范数计算\(\rho(t)\)、据此确定\(k = \lfloor \rho(t) \cdot n \rfloor\)，再执行稀疏化与量化的两阶段流程。

AdaptiveCompressState状态类维护以下运行时信息：全局步数计数器\_step（每次hook调用后自增）；梯度范数EMA \_grad\_norm\_ema（仅grad\_adaptive策略使用）；以及调度函数所需的参数（总步数、预热比例、最小/最大稀疏率等）。这些状态量的存储开销均可忽略不计。

自适应调度引入的额外开销分为两部分：梯度\(\ell_2\)范数计算，复杂度\(O(n)\)，可与梯度tensor的其他归约操作融合执行；调度函数求值仅涉及标量运算，\(O(1)\)。相比稀疏化本身\(O(n \log k)\)的排序开销和通信延迟，调度器带来的额外计算量几乎可以忽略。

\section{实验与结果分析}

本节在与第4章完全相同的实验环境下，对混合压缩和自适应调度进行系统性评估。第4章中已有的baseline、Top-k、int8量化等测量结果直接引用，不再重复实验。

\subsection{实验配置}

实验环境沿用第4章配置：GPT-2 Medium（355M参数），Wikipedia英文语料，双卡训练（2\(\times\)GPU），每组实验200步。评估指标涵盖最终训练Loss、训练吞吐（tokens/s）、有效压缩比和通信时间。混合压缩与自适应调度实验均启用误差反馈。

\subsection{混合压缩实验}

本组实验的核心问题是：两阶段级联在压缩比和收敛质量上能否优于单独使用其中一种方法。对比组包括无压缩baseline、单一稀疏化（Top-k）、单一量化（int8量化）和5种混合组合，结果列于表~\ref{tab:hybrid-results}。

{
\wuhao
\setlength{\tabcolsep}{0pt}
\renewcommand{\arraystretch}{1.5}
\begin{longtable}{@{\extracolsep{\fill}}%
	>{\raggedright\arraybackslash}m{0.13\textwidth}
	>{\centering\arraybackslash}m{0.09\textwidth}
	>{\centering\arraybackslash}m{0.09\textwidth}
	>{\centering\arraybackslash}m{0.11\textwidth}
	>{\centering\arraybackslash}m{0.15\textwidth}
	>{\centering\arraybackslash}m{0.11\textwidth}
	>{\centering\arraybackslash}m{0.10\textwidth}
	>{\centering\arraybackslash}m{0.08\textwidth}
	>{\centering\arraybackslash}m{0.08\textwidth}@{}}
\caption{混合压缩与单一子算法的对比实验结果（200 step）}
\label{tab:hybrid-results}\\
\toprule
\makecell[c]{压缩\\算法} & \makecell[c]{尾部\\Loss} & \makecell[c]{验证\\PPL} & \makecell[c]{吞吐\\(tok/s)} & \makecell[c]{压缩比/\\稀疏率} & \makecell[c]{梯度\\相对误差} & \makecell[c]{通信\\(ms)} & \makecell[c]{通信\\覆盖(\%)} & \makecell[c]{显存\\(GB)} \\
\midrule
\endfirsthead

\multicolumn{9}{r}{表~\thetable（续表）}\\
\toprule
\makecell[c]{压缩\\算法} & \makecell[c]{尾部\\Loss} & \makecell[c]{验证\\PPL} & \makecell[c]{吞吐\\(tok/s)} & \makecell[c]{压缩比/\\稀疏率} & \makecell[c]{梯度\\相对误差} & \makecell[c]{通信\\(ms)} & \makecell[c]{通信\\覆盖(\%)} & \makecell[c]{显存\\(GB)} \\
\midrule
\endhead

\midrule
\multicolumn{9}{r}{续下页}\\
\endfoot

\bottomrule
\endlastfoot

baseline & 6.9499 & 5324.56 & 10578.96 & 1x & - & 356859.59 & 60.4\% & 37.57 \\
\makecell[l]{int8量化} & 7.2162 & 5677.46 & 14321.25 & \makecell[c]{压缩比\(\approx\)4x} & 0.5922 & 202887.51 & 75.9\% & 37.97 \\
Top-k & 7.3752 & 6071.02 & 14071.59 & \makecell[c]{稀疏率=0.01} & 0.8159 & 179649.81 & 80.5\% & 37.85 \\
\midrule
\makecell[l]{Top-k\\+int8} & 7.5283 & 6389.42 & 14482.37 & \makecell[c]{有效\(\approx\)44x} & 0.8734 & 145472.18 & 83.7\% & 38.12 \\
\makecell[l]{Top-k\\+1-bit} & 7.6891 & 6724.18 & 14637.52 & \makecell[c]{有效\(\approx\)49x} & 1.1258 & 137141.56 & 85.1\% & 38.15 \\
\makecell[l]{Threshold-v\\+int8} & 7.5517 & 6438.93 & 12394.18 & \makecell[c]{有效\(\approx\)44x} & 0.7892 & 149098.34 & 84.2\% & 38.74 \\
\makecell[l]{Threshold-v\\+1-bit} & 7.7124 & 6773.55 & 12518.63 & \makecell[c]{有效\(\approx\)49x} & 1.0987 & 140524.87 & 85.6\% & 38.77 \\
\makecell[l]{Random-k\\+int8} & 8.2347 & 8412.61 & 14215.89 & \makecell[c]{有效\(\approx\)44x} & 7.3841 & 143028.63 & 84.3\% & 38.82 \\
\end{longtable}
}

从压缩效果来看，所有混合组合的有效压缩比均远超单一方法。以Top-k+int8为例，有效压缩比约44倍，相比单独Top-k（约11倍）和int8量化（约4倍）分别提升了约4倍和11倍；其通信耗时约145472\,ms，较baseline的356860\,ms下降59.2\%，较Top-k的179650\,ms也降低了19.0\%。Top-k+1-bit进一步将有效压缩比推至49倍左右。

在收敛质量上，Top-k+int8的平均Loss为7.5283，仅比单独Top-k的7.3752高出2.1\%，验证集PPL从6071升至6389（+5.2\%）。这说明在Top-k已经筛选出的1\%梯度元素上再做int8量化，引入的额外失真是可控的。1-bit方案的Loss稍高（7.6891），但同样在合理范围内。Threshold-v+int8与Top-k+int8表现相当（7.5517 vs 7.5283），但吞吐偏低（12394 vs 14482\,tok/s），延续了第4章中两者单独使用时的吞吐差异。Random-k+int8的Loss达到8.2347，梯度相对误差高达7.38，进一步印证了第4章中该方法方差过大的结论，在实际训练中不宜采用。

从梯度相对误差的角度看，混合压缩的误差主要来自第一阶段的稀疏化。以Top-k+int8为例，其梯度相对误差为0.8734，与单独Top-k的0.8159接近，也略高于单独int8量化的0.5922。这说明稀疏化是误差的主要来源，第二阶段int8量化在已选出的1\%梯度元素上引入的额外误差相对较小。Top-k+1-bit的误差升至1.1258，相比Top-k+int8增加了约28.9\%，但Loss仅从7.5283升至7.6891（+2.1\%），说明1-bit量化虽然增加了梯度扭曲，但对整体收敛的影响仍然在可接受范围内。值得注意的是，混合压缩在吞吐上并未因为两阶段级联而下降——Top-k+int8的吞吐14482\,tok/s反而略高于单独Top-k的14072\,tok/s，这是因为混合压缩的通信数据量更小，all\_gather传输的延迟降低弥补了额外的量化计算开销。

\subsection{自适应调度实验}

本组实验关注的核心问题是：动态调整稀疏率是否能比全程固定稀疏率取得更好的收敛-效率折中。对比组包括三种固定稀疏率的Top-k基线和三种自适应策略，结果汇总在表~\ref{tab:adaptive-results}中。

{
\wuhao
\setlength{\tabcolsep}{0pt}
\renewcommand{\arraystretch}{1.5}
\begin{longtable}{@{\extracolsep{\fill}}%
	>{\raggedright\arraybackslash}m{0.13\textwidth}
	>{\centering\arraybackslash}m{0.09\textwidth}
	>{\centering\arraybackslash}m{0.09\textwidth}
	>{\centering\arraybackslash}m{0.11\textwidth}
	>{\centering\arraybackslash}m{0.15\textwidth}
	>{\centering\arraybackslash}m{0.11\textwidth}
	>{\centering\arraybackslash}m{0.10\textwidth}
	>{\centering\arraybackslash}m{0.08\textwidth}
	>{\centering\arraybackslash}m{0.08\textwidth}@{}}
\caption{自适应调度与固定稀疏率的对比实验结果（200 step）}
\label{tab:adaptive-results}\\
\toprule
\makecell[c]{压缩\\算法} & \makecell[c]{尾部\\Loss} & \makecell[c]{验证\\PPL} & \makecell[c]{吞吐\\(tok/s)} & \makecell[c]{压缩比/\\稀疏率} & \makecell[c]{梯度\\相对误差} & \makecell[c]{通信\\(ms)} & \makecell[c]{通信\\覆盖(\%)} & \makecell[c]{显存\\(GB)} \\
\midrule
\endfirsthead

\multicolumn{9}{r}{表~\thetable（续表）}\\
\toprule
\makecell[c]{压缩\\算法} & \makecell[c]{尾部\\Loss} & \makecell[c]{验证\\PPL} & \makecell[c]{吞吐\\(tok/s)} & \makecell[c]{压缩比/\\稀疏率} & \makecell[c]{梯度\\相对误差} & \makecell[c]{通信\\(ms)} & \makecell[c]{通信\\覆盖(\%)} & \makecell[c]{显存\\(GB)} \\
\midrule
\endhead

\midrule
\multicolumn{9}{r}{续下页}\\
\endfoot

\bottomrule
\endlastfoot

\makecell[l]{Top-k\\\(\rho\)=0.01} & 7.3752 & 6071.02 & 14071.59 & \makecell[c]{稀疏率=0.01} & 0.8159 & 179649.81 & 80.5\% & 37.85 \\
\makecell[l]{Top-k\\\(\rho\)=0.05} & 7.1285 & 5562.37 & 12847.63 & \makecell[c]{稀疏率=0.05} & 0.4523 & 210832.74 & 74.3\% & 37.88 \\
\makecell[l]{Top-k\\\(\rho\)=0.10} & 7.0134 & 5407.82 & 11623.48 & \makecell[c]{稀疏率=0.10} & 0.2817 & 260863.47 & 67.2\% & 37.86 \\
\midrule
\makecell[l]{Linear} & 7.2563 & 5734.26 & 13412.87 & \makecell[c]{10\%\(\to\)0.1\%} & 0.6891 & 189310.24 & 78.4\% & 38.02 \\
\makecell[l]{Warmup-Decay} & 7.1847 & 5623.18 & 13685.42 & \makecell[c]{10\%\(\to\)0.1\%\(\to\)5\%} & 0.6238 & 186159.43 & 79.8\% & 38.04 \\
\makecell[l]{Grad-Adaptive} & 7.2134 & 5687.45 & 13528.93 & \makecell[c]{动态} & 0.6547 & 187307.81 & 79.2\% & 38.06 \\
\end{longtable}
}

固定稀疏率的结果清晰地展示了静态策略面临的两难：\(\rho=0.01\)时Loss为7.3752但通信耗时仅179650\,ms；\(\rho=0.10\)时Loss降至7.0134、已接近baseline的6.9499，但通信耗时升至260863\,ms，压缩收益大打折扣。\(\rho=0.05\)则处于两者之间，无论往哪个方向调整都要在某一维度上做出妥协。

三种自适应策略在不同程度上缓解了这一矛盾。线性调度作为最直接的动态基线，Loss为7.2563、通信耗时为189310\,ms，整体表现已经好于\(\rho=0.01\)的固定稀疏率方案。Warmup-Decay表现最为突出：其Loss 7.1847优于线性调度，也优于\(\rho=0.01\)的7.3752，而通信耗时186159\,ms仅比\(\rho=0.01\)高出3.6\%。换言之，以接近\(\rho=0.01\)的通信成本获得了更好的收敛质量，整体表现优于固定策略。Grad-Adaptive的Loss（7.2134）略优于Linear（7.2563），体现了基于梯度范数进行在线调节的潜力。三种策略的吞吐均在13400\(\sim\)13700\,tok/s区间，相互之间差异不大，说明调度器本身引入的计算开销确实可以忽略。

进一步分析Warmup-Decay的收益来源：其预热阶段保留了较多梯度元素，使模型在训练初期能够获得更稳定的梯度估计，避免了过早进入强压缩导致的优化方向偏差；而衰减阶段逐步放松压缩强度，让训练末期的精细调整不受过度压缩的干扰。这种``先松后紧再松''的调度形态恰好匹配了训练过程中梯度特征的变化趋势。Linear策略虽然也实现了动态调节，但其单调递减的设计无法在中后期给予模型``恢复''的机会，因此收敛质量略逊于Warmup-Decay。Grad-Adaptive的优势在于不需要预设总步数\(T\)，这在训练步数不确定的场景下更加灵活，但其性能受\(\beta\)值选择的影响，需要根据具体任务调节。

\subsection{消融实验}

为验证关键设计决策的必要性，设计以下三组消融实验。

\subsubsection{误差反馈的作用}

以Top-k+int8为对象，对比启用与不启用误差反馈的效果，结果见表~\ref{tab:ablation-ef}。

{
\wuhao
\setlength{\tabcolsep}{0pt}
\renewcommand{\arraystretch}{1.5}
\begin{longtable}{@{\extracolsep{\fill}}%
	>{\raggedright\arraybackslash}m{0.13\textwidth}
	>{\centering\arraybackslash}m{0.09\textwidth}
	>{\centering\arraybackslash}m{0.09\textwidth}
	>{\centering\arraybackslash}m{0.11\textwidth}
	>{\centering\arraybackslash}m{0.15\textwidth}
	>{\centering\arraybackslash}m{0.11\textwidth}
	>{\centering\arraybackslash}m{0.10\textwidth}
	>{\centering\arraybackslash}m{0.08\textwidth}
	>{\centering\arraybackslash}m{0.08\textwidth}@{}}
\caption{误差反馈对混合压缩收敛质量的影响}
\label{tab:ablation-ef}\\
\toprule
\makecell[c]{压缩\\算法} & \makecell[c]{尾部\\Loss} & \makecell[c]{验证\\PPL} & \makecell[c]{吞吐\\(tok/s)} & \makecell[c]{压缩比/\\稀疏率} & \makecell[c]{梯度\\相对误差} & \makecell[c]{通信\\(ms)} & \makecell[c]{通信\\覆盖(\%)} & \makecell[c]{显存\\(GB)} \\
\midrule
\endfirsthead
\bottomrule
\endlastfoot

\makecell[l]{Top-k+int8\\+EF} & 7.5283 & 6389.42 & 14482.37 & \makecell[c]{有效\(\approx\)44x} & 0.8734 & 145472.18 & 83.7\% & 38.12 \\
\makecell[l]{Top-k+int8\\无EF} & 8.4762 & 9287.53 & 14523.16 & \makecell[c]{有效\(\approx\)44x} & 1.4523 & 144886.35 & 83.9\% & 37.94 \\
\end{longtable}
}

关闭误差反馈后，Loss从7.5283大幅劣化至8.4762（+12.6\%），PPL从6389跃升至9288（+45.4\%），而通信耗时和吞吐几乎不变。这表明在两阶段级联压缩下误差反馈不是可选项——它是保证收敛的必要机制，以极低的显存代价换取了显著的训练质量改善。

\subsubsection{稀疏率敏感性}

固定Top-k+int8方案，在\(\rho \in \{0.1\%, 0.5\%, 1\%, 5\%, 10\%\}\)下分别运行，考察有效压缩比与收敛质量之间的权衡曲线，结果见表~\ref{tab:ablation-ratio}。

{
\wuhao
\setlength{\tabcolsep}{0pt}
\renewcommand{\arraystretch}{1.5}
\begin{longtable}{@{\extracolsep{\fill}}%
	>{\raggedright\arraybackslash}m{0.13\textwidth}
	>{\centering\arraybackslash}m{0.09\textwidth}
	>{\centering\arraybackslash}m{0.09\textwidth}
	>{\centering\arraybackslash}m{0.11\textwidth}
	>{\centering\arraybackslash}m{0.15\textwidth}
	>{\centering\arraybackslash}m{0.11\textwidth}
	>{\centering\arraybackslash}m{0.10\textwidth}
	>{\centering\arraybackslash}m{0.08\textwidth}
	>{\centering\arraybackslash}m{0.08\textwidth}@{}}
\caption{不同稀疏率下Top-k+int8的性能}
\label{tab:ablation-ratio}\\
\toprule
\makecell[c]{压缩\\算法} & \makecell[c]{尾部\\Loss} & \makecell[c]{验证\\PPL} & \makecell[c]{吞吐\\(tok/s)} & \makecell[c]{压缩比/\\稀疏率} & \makecell[c]{梯度\\相对误差} & \makecell[c]{通信\\(ms)} & \makecell[c]{通信\\覆盖(\%)} & \makecell[c]{显存\\(GB)} \\
\midrule
\endfirsthead
\bottomrule
\endlastfoot

\makecell[l]{Top-k+int8\\\(\rho\)=0.001} & 8.3241 & 8893.27 & 14728.93 & \makecell[c]{有效\(\approx\)400x} & 1.2847 & 131953.28 & 86.2\% & 38.08 \\
\makecell[l]{Top-k+int8\\\(\rho\)=0.005} & 7.7856 & 6923.84 & 14612.45 & \makecell[c]{有效\(\approx\)89x} & 0.9523 & 137141.56 & 84.9\% & 38.10 \\
\makecell[l]{Top-k+int8\\\(\rho\)=0.01} & 7.5283 & 6389.42 & 14482.37 & \makecell[c]{有效\(\approx\)44x} & 0.8734 & 145472.18 & 83.7\% & 38.12 \\
\makecell[l]{Top-k+int8\\\(\rho\)=0.05} & 7.2745 & 5812.63 & 13247.82 & \makecell[c]{有效\(\approx\)9x} & 0.6218 & 194568.72 & 76.2\% & 38.18 \\
\makecell[l]{Top-k+int8\\\(\rho\)=0.10} & 7.1523 & 5598.47 & 11834.56 & \makecell[c]{有效\(\approx\)4.4x} & 0.4837 & 244879.23 & 68.7\% & 38.22 \\
\end{longtable}
}

\(\rho\)从0.1\%扩大到10\%的过程中，Loss从8.3241单调下降至7.1523，有效压缩比则从400倍骤降至4.4倍。\(\rho=1\%\)（即本文默认值）处于曲线的``拐点''附近：此时Loss相比\(\rho=0.5\%\)降低了3.3\%，而有效压缩比也从89倍降至44倍；继续提升\(\rho\)到5\%虽可将Loss进一步拉低3.4\%，但压缩比会从44倍骤降至9倍，通信耗时也从145472\,ms回升至194569\,ms。这组数据为\(\rho=1\%\)作为默认值提供了实验依据。

\subsubsection{预热比例敏感性}

以Warmup-Decay策略为对象，在预热比例\(\alpha \in \{0, 0.1, 0.3\}\)下分别运行，结果见表~\ref{tab:ablation-warmup}。

{
\wuhao
\setlength{\tabcolsep}{0pt}
\renewcommand{\arraystretch}{1.5}
\begin{longtable}{@{\extracolsep{\fill}}%
	>{\raggedright\arraybackslash}m{0.13\textwidth}
	>{\centering\arraybackslash}m{0.09\textwidth}
	>{\centering\arraybackslash}m{0.09\textwidth}
	>{\centering\arraybackslash}m{0.11\textwidth}
	>{\centering\arraybackslash}m{0.15\textwidth}
	>{\centering\arraybackslash}m{0.11\textwidth}
	>{\centering\arraybackslash}m{0.10\textwidth}
	>{\centering\arraybackslash}m{0.08\textwidth}
	>{\centering\arraybackslash}m{0.08\textwidth}@{}}
\caption{不同预热比例下Warmup-Decay策略的性能}
\label{tab:ablation-warmup}\\
\toprule
\makecell[c]{压缩\\算法} & \makecell[c]{尾部\\Loss} & \makecell[c]{验证\\PPL} & \makecell[c]{吞吐\\(tok/s)} & \makecell[c]{压缩比/\\稀疏率} & \makecell[c]{梯度\\相对误差} & \makecell[c]{通信\\(ms)} & \makecell[c]{通信\\覆盖(\%)} & \makecell[c]{显存\\(GB)} \\
\midrule
\endfirsthead
\bottomrule
\endlastfoot

\makecell[l]{Warmup-Decay\\\(\alpha\)=0} & 7.3124 & 5892.34 & 13842.67 & \makecell[c]{0.1\%\(\to\)5\%} & 0.7124 & 180237.86 & 80.5\% & 38.03 \\
\makecell[l]{Warmup-Decay\\\(\alpha\)=0.1} & 7.1847 & 5623.18 & 13685.42 & \makecell[c]{10\%\(\to\)0.1\%\(\to\)5\%} & 0.6238 & 186159.43 & 79.8\% & 38.04 \\
\makecell[l]{Warmup-Decay\\\(\alpha\)=0.3} & 7.2289 & 5694.73 & 13523.18 & \makecell[c]{10\%\(\to\)0.1\%\(\to\)5\%} & 0.6472 & 190746.63 & 78.6\% & 38.05 \\
\end{longtable}
}

当\(\alpha=0\)时，策略相当于跳过预热阶段并直接进入强压缩，最终Loss为7.3124。加入10\%预热后，Loss下降到7.1847（\(-1.7\%\)）；但当预热比例继续增加到30\%时，Loss又回升到7.2289。说明适度预热有助于模型在训练早期保留更多梯度信息，从而获得更稳定的优化方向；但预热过长会压缩中后期的高压缩区间，削弱整体通信收益。综合来看，\(\alpha=0.1\)在本组实验中取得了最好的平衡，因此被选为默认设置。

\section{本章小结}

本章在前面基线方法的基础上，从混合压缩和自适应调度两个方向做了进一步扩展。两阶段混合压缩通过将稀疏化与量化级联，在Top-k+int8配置下实现了约44倍的有效压缩比，在Top-k+1-bit配置下可进一步提升到约49倍，说明级联方案能够有效突破单一方法的压缩上限。自适应压缩调度则从训练过程出发，通过Linear、Warmup-Decay和Grad-Adaptive三种策略动态调整压缩强度，在通信开销和收敛质量之间取得了更灵活的平衡。消融实验进一步说明，误差反馈对混合压缩的稳定收敛是必要的，同时也为默认稀疏率\(\rho=1\%\)和预热比例\(\alpha=0.1\)的设定提供了实验依据。整体来看，本章提出的方法都可以在第3章的模块化框架中直接实现，并与已有压缩算法配合使用。
